\documentclass[12pt,letterpaper]{article}
\usepackage[T1]{fontenc}
\usepackage[utf8]{inputenc}
\usepackage{lmodern}
\usepackage[margin=1in]{geometry}
\usepackage{amsmath,amssymb,amsthm,mathtools,mathrsfs}
\usepackage{microtype,booktabs,array}
\usepackage[hidelinks]{hyperref}
\usepackage{bookmark}
\setlength{\emergencystretch}{2em}
\newtheorem{theorem}{Theorem}[section]
\newtheorem{lemma}[theorem]{Lemma}
\newtheorem{proposition}[theorem]{Proposition}
\newtheorem{corollary}[theorem]{Corollary}
\theoremstyle{definition}\newtheorem{definition}[theorem]{Definition}
\theoremstyle{remark}\newtheorem{remark}[theorem]{Remark}
\DeclareMathOperator{\Tr}{Tr}
\DeclareMathOperator{\rank}{rank}
\DeclareMathOperator{\Var}{Var}
\DeclareMathOperator{\supp}{supp}
\newcommand{\C}{\mathbb C}
\newcommand{\N}{\mathbb N}
\newcommand{\E}{\mathbb E}
\newcommand{\Prb}{\mathbb P}
\newcommand{\id}{\mathrm{id}}
\newcommand{\LO}{\mathrm{LO}}
\newcommand{\LOCC}{\mathrm{LOCC}}
\newcommand{\cPPT}{\mathrm{cPPT}}
\newcommand{\APPT}{\mathrm{APPT}}
\newcommand{\G}{\mathcal G}
\newcommand{\calC}{\mathcal C}
\newcommand{\pos}[1]{(#1)_+}
\newcommand{\ip}[2]{\langle #1,#2\rangle}
\hypersetup{pdftitle={The finite-block cost of universality in entanglement extraction},pdfauthor={Yongxian Zhang},pdfkeywords={universal quantum protocols, entanglement extraction, finite blocklength, minimax, unknown support}}
\title{The finite-block cost of universality\\in entanglement extraction}
\author{Yongxian Zhang\thanks{School of Computer Science and Engineering, South China University of Technology, Guangzhou, China. Email: \href{mailto:mxymmxym1@gmail.com}{mxymmxym1@gmail.com}. ORCID: \href{https://orcid.org/0009-0000-3864-3536}{0009-0000-3864-3536}.}}
\date{October 2026}
\begin{document}
\maketitle
\begin{abstract}
We study entanglement extraction by a global unitary on an existing many-copy input followed by deterministic local channels. The global preprocessing is an explicit entangling resource; no auxiliary system is supplied before it. We establish finite-dimensional spectral and subspace packing principles that match converses for completely PPT postprocessing. One state-independent protocol attains the state-aware capacity, exact fidelity exponent, strict-direct reliability function, and entropy-limited Gaussian threshold of every fixed input. A finite-block information-spectrum sandwich retains its dimension dependence and gives sufficient conditions for growing-dimensional Gaussian limits. We then determine the sharp minimax cost of an unknown support. For every fixed rank budget, a coherent common-support code matches a Haar-average converse, relatively and uniformly over all target dimensions. The support dimension and the reciprocal purity of the averaged input have the same leading asymptotic, yielding an explicit logarithmic loss in the extractable target size and a transition at the local-dimension cap. For balanced unknown pure inputs, the full-output fidelity has leading constant one. For unknown pure inputs, an explicit original-register compiler outputs finite-library circuits with a proved word-length and uniform fidelity guarantee. These results distinguish pointwise asymptotic universality from uniform finite-block performance. Standard compression, representation theory, and distillation results are explicitly credited.
\end{abstract}

\section{Introduction and operational model}

Universal quantum protocols operate without a classical description of their input. Universal compression, pure-state concentration, and mixed-state distillation establish several senses in which ignorance need not change an asymptotic rate \cite{jhhh,hm,takagi,tirone,llf,rizzo}. Equality of rates, however, does not determine a common protocol's finite-block performance. The same leading capacity can coexist with a polynomial fidelity loss, and a rank promise can incur a logarithmic target-size penalty invisible to the first-order rate.

We examine these distinctions in a two-stage extraction model. First apply one global unitary on the existing input space. Then apply deterministic local product channels, LOCC, or completely PPT channels. The global unitary is an entangling resource. This task is therefore distinct from ordinary LOCC distillation of the original fixed state. Its utility is to separate the spectral supply of purity from subsequent extraction and to ask what can be achieved by one preprocessing independent of the entire state.

The main method is a finite-dimensional comparison between spectral or common-subspace packing and the quadratic constraint on a completely PPT fidelity effect. It has three consequences. First, a variable-size Bell packing gives a one-shot spectral envelope, within logarithmic factors, for every bipartite spectrum. Second, intrinsic polynomial matrix-orbit projections and a geometric Bell-packing grid give one protocol for all states, with an explicit finite-block information-spectrum sandwich. Third, for a class of unknown inputs, a common-support code is compared with an adversarial averaged state. Literal Haar minor moments and an operator lower bound show that the common-support dimension and the reciprocal purity of that average agree asymptotically. The entire minimax fidelity profile then follows with leading constant one, uniformly over all target integers.

\subsection{Model and quantifiers}
Fix integers $2\le m\le n$, set $D=mn$, and write
\[
 a=m^k,\qquad b=n^k,\qquad N=D^k
\]
at copy number $k\ge1$. All logarithms are natural. Let $\Phi_K$ denote the normalized rank-one maximally entangled projector on $\C^K\otimes\C^K$. We use the target overlap $\Tr(\Phi_K\sigma)$ as fidelity; no square root is taken.

A protocol comprises a unitary $U$ on $\C^a\otimes\C^b$ and a trace-preserving channel $\Lambda$ to $\C^K\otimes\C^K$. No ancilla, measurement, or discarded register is supplied before $U$. Local ancillas and discarding are permitted afterwards. The operation class $\LO$ consists of product channels, and $\LOCC$ denotes finite-round adaptive local instruments with communicated outcomes; hidden Kraus indices may refine each public outcome. We also consider its closure in the finite-dimensional topology of all matrix-unit channel coefficients. A channel is completely PPT, denoted $\cPPT$, if both $\Lambda$ and $\Gamma_{\rm out}\Lambda\Gamma_{\rm in}$ are completely positive and trace preserving; $\Gamma$ transposes the first party. Thus $\LO\subseteq\LOCC\subseteq\overline{\LOCC}\subseteq\cPPT$. All stated intermediate-class conclusions apply to the closure as well. Complete positivity of the transpose-conjugated channel is required on every auxiliary space. The arbitrary global preprocessing $U$ remains a separate entangling resource and is not included in this PPT restriction.

For a density matrix $\rho$ on $\C^m\otimes\C^n$, define
\begin{equation}\label{eq:state-aware}
 f_k^{\calC}(\rho,K)=\sup_{U,\Lambda\in\calC}
 \Tr\!\left[\Phi_K\Lambda(U\rho^{\otimes k}U^*)\right].
\end{equation}
The optimizing protocol in this expression may depend on $\rho$. In contrast, our all-state universal family has quantifiers
\begin{equation}\label{eq:univ-quantifiers}
 \exists\{U_{k,K},\mathcal A_{k,K},\mathcal B_{k,K}\}_{k,K}
 \quad\forall\rho,
\end{equation}
with dependence only on $m,n,k,K$. We write $F_k^{\rm u}(\rho,K)$ for its fidelity and $F_k^\star=f_k^{\cPPT}$. No spectrum, eigenbasis, rank, or entropy is an input to this family.

A different, explicitly promised task fixes a rank budget $1\le s\le D$ and defines
\begin{equation}\label{eq:minimax-definition}
 \G_{\calC}(k,K;s)=\sup_{U,\Lambda\in\calC}
 \inf_{\rank\rho\le s}\Tr\!\left[\Phi_K\Lambda(U\rho^{\otimes k}U^*)\right].
\end{equation}
Here one protocol may depend on $s$, but it must serve every spectrum and every support orientation in that class. Pointwise universality in \eqref{eq:univ-quantifiers} and minimax universality in \eqref{eq:minimax-definition} are different assertions. We retain that distinction throughout.

\subsection{Results and relation to prior work}
Section~\ref{sec:spectral} proves the actual finite-dimensional converse and packing principles. Section~\ref{sec:universal} gives a universal physical construction and a dimension-explicit finite-block sandwich. Section~\ref{sec:asymptotic} derives the exact capacity and error laws. Section~\ref{sec:support} proves the sharp rank-budget minimax profile, including the boundary at $n=ms$ and its logarithmic correction. The pure-promise construction also has an explicit compiler on the original register, with a fixed finite gate library and quantitative whole-circuit error control. Section~\ref{sec:applications} connects the capacity to APPT entropy optimization and specifies the scope of the accompanying formal verification.

The fidelity-effect characterization is due to Rains \cite{rains}. Lami's quadratic converse already implies the half entropy-deficit upper bound \cite{lami}; we supply a direct finite-dimensional derivation and matching local constructions. Schur--Weyl universal compression and estimation are established methods \cite{jhhh,kw,hayashi-univ}. Exact symmetric-state compression, including efficient qubit constructions, predates this work \cite{plesch,baertschi}. We do not claim the symmetric-subspace dimension or compression mechanism as new. Classical source-coding reliability and second-order asymptotics are likewise prior results \cite{hayashi-second}.

Recent universal distillation and coding results remove input-state knowledge under other free operations \cite{takagi,tirone,llf,rizzo}. State-agnostic work extraction studies a thermodynamic model \cite{watanabe,zeng}. Their operations and resources differ from the global entangling preprocessing allowed here. Our claims concern the stated two-stage model: a common deterministic physical packing, matching completely PPT converses, finite-block control with visible dimension dependence, and a sharp target-uniform minimax law for an unknown rank-bounded support. The fact that a protocol is universal, by itself, is not presented as a new principle.

Known-spectrum, unknown-eigenbasis families also appear in the quantum minimum-description-length problem of Hayden, Maloney, Wang and Yang \cite{hayden-qmdl}. Their encoder and decoder minimize the memory needed to recover the many-copy state with vanishing trace-distance error, without requiring preservation of its purification. The loss quantified here is the attainable entangled target under a common global unitary and deterministic local postprocessing, with a uniform relative fidelity profile over all target dimensions.

Purity-based extraction also includes the stochastic-unital analysis of Kondra et al.\ \cite{kondra}, which characterizes complete absolute separability, studies stability under adjoining maximally mixed local systems, and proves finite-copy activation. Their probabilistic convertibility problem permits postselection and compares states in the same output space. Our deterministic, variable-target fidelity optimization has a different resource specification. The qualitative possibility of activation is prior work.

\section{Finite-dimensional spectral and subspace principles}\label{sec:spectral}

Throughout this section the input dimensions are arbitrary integers $2\le a\le b$, $N=ab$, and $K\ge2$. The case $K=1$ is the trivial exact product target.

\begin{lemma}[Fidelity effects, following Rains]\label{lem:rains}
For a completely PPT channel let $M=\Lambda^*(\Phi_K)$. Then
\begin{equation}\label{eq:rains}
 0\le M\le I,\qquad -I/K\le M^\Gamma\le I/K.
\end{equation}
Conversely every such $M$ is a fidelity effect of a completely PPT channel. Moreover,
\begin{equation}\label{eq:budget}
 \Tr M^2\le N/K^2,\qquad
 \|M\|_\infty\le\beta:=\min(1,a/K).
\end{equation}
\end{lemma}
\begin{proof}
The adjoint of a channel is positive and unital. Since $\Phi_K^\Gamma=F_K/K$ for the swap $F_K$, and $-I\le F_K\le I$, apply the adjoint of $\Gamma_{\rm out}\Lambda\Gamma_{\rm in}$ to these order bounds. This gives \eqref{eq:rains}.

For the converse define
\[
 \Lambda_M(X)=\Tr(MX)\Phi_K+
 \Tr((I-M)X)\frac{I-\Phi_K}{K^2-1}.
\]
This is a measure-and-prepare channel. Writing $Q_\pm=(I\pm F_K)/2$, its partial-transpose conjugate is
\[
 X\longmapsto
 \Tr\!\left[\frac{I+KM^\Gamma}{K(K+1)}X\right]Q_+
 +\Tr\!\left[\frac{I-KM^\Gamma}{K(K-1)}X\right]Q_-.
\]
Both coefficient effects are positive, and the map is trace preserving because $\Tr Q_\pm=K(K\pm1)/2$. Its target overlap is $\Tr(MX)$.

Partial transpose preserves the Hilbert--Schmidt norm. The eigenvalues of $M^\Gamma$ belong to $[-1/K,1/K]$, proving the first bound in \eqref{eq:budget}. For a unit vector $v$, the partial transpose of $vv^*$ has trace norm $(\sum_j t_j)^2\le a$, where $t_j$ are the Schmidt coefficients. Hence
\[
 \langle v,Mv\rangle
 =\Tr\bigl(M^\Gamma(vv^*)^\Gamma\bigr)\le a/K.
\]
Together with $M\le I$, this proves the norm bound.
\end{proof}

\begin{lemma}[Coherent common-subspace packing]\label{lem:packing}
Let $S\subseteq\C^a\otimes\C^b$ be a known subspace of dimension $r$. If
\begin{equation}\label{eq:packing-feasible}
 1\le d\le\min(a,K),\qquad
 \lfloor a/d\rfloor\lfloor b/d\rfloor\ge r,
\end{equation}
one unitary on the original input followed by deterministic product channels sends every density matrix supported on $S$ to the same embedded $\Phi_d$. Its target-$K$ fidelity is $d/K$.
\end{lemma}
\begin{proof}
Put $u=\lfloor a/d\rfloor$ and $v=\lfloor b/d\rfloor$. In local coordinate blocks there are $uv$ orthonormal vectors
\[
 |\psi_{ij}\rangle=d^{-1/2}\sum_{t=0}^{d-1}
 |id+t\rangle_A|jd+t\rangle_B,
 \quad 0\le i<u,\quad0\le j<v.
\]
Map a fixed orthonormal basis of $S$ to $r$ distinct such vectors and extend the resulting isometry to a unitary on the full $N$-dimensional input. Alice uses
\[
 A_i=\sum_{t<d}|t\rangle_K\langle id+t|,
\]
and $|0\rangle\langle z|$ for each unused input coordinate $z$. Bob uses the corresponding maps. The Kraus adjoint products sum to $I_a$ and $I_b$, including all unused coordinates.

For these product channels,
\[
 |\psi_{ij}\rangle\langle\psi_{i'j'}|
 \longmapsto \delta_{ii'}\delta_{jj'}\Phi_d.
\]
Thus any logical matrix $X$ on the selected code maps to $(\Tr X)\Phi_d$, including every off-diagonal coherence. A density matrix on $S$ therefore has output $\Phi_d$. Finally $\Tr(\Phi_K\Phi_d)=d/K$. No branch is postselected or renormalized.
\end{proof}

\begin{proposition}[A subspace-versus-average comparison]\label{prop:average}
Let a family $\mathscr F$ of input states be supported on one $r$-dimensional subspace, and let $\Omega$ be a probability average of states in that family. Define
\[
 d_{a,b}(r)=\max\{d\in\{1,\ldots,a\}:\lfloor a/d\rfloor\lfloor b/d\rfloor\ge r\}.
\]
For the minimax fidelity over $\mathscr F$,
\begin{align}\label{eq:general-minimax}
 \min(1,d_{a,b}(r)/K)
 &\le\sup_{U,\Lambda\in\LO}\inf_{\sigma\in\mathscr F}F(U,\Lambda;\sigma)\nonumber\\
 &\le\sup_{U,\Lambda\in\cPPT}\inf_{\sigma\in\mathscr F}F(U,\Lambda;\sigma)
 \le\min\left(1,\frac aK,\frac{\sqrt{ab\Tr\Omega^2}}K\right).
\end{align}
The intermediate LOCC minimax value satisfies the same bounds.
\end{proposition}
\begin{proof}
Use Lemma~\ref{lem:packing} with $d=\min(K,d_{a,b}(r))$. For the upper bound, the infimum for a fixed protocol is at most its average,
\[
 \inf_{\sigma\in\mathscr F}F(U,\Lambda;\sigma)
 \le\Tr(MU\Omega U^*)
 \le\sqrt{\Tr M^2\Tr\Omega^2}.
\]
Apply \eqref{eq:budget}; the other two bounds follow from the trace-one condition and the norm bound. Averaging is used only to prove a converse, not as an operation in the protocol.
\end{proof}

We also need a spectral comparison for a single input. Let its eigenvalues be $\lambda_1\ge\cdots\ge\lambda_N\ge0$, and put
\begin{equation}\label{eq:envelope}
 L_r=\sum_{i\le r}\lambda_i,\quad
 v_r=\min\left(1,\frac aK,\frac{\sqrt{N/r}}K\right),\quad
 T_K=\max_{1\le r\le N}v_rL_r.
\end{equation}

\begin{theorem}[One-shot spectral envelope]\label{thm:envelope}
For every input spectrum and every $K\ge2$,
\begin{equation}\label{eq:envelope-bound}
 T_K/4\le f^{\LO}(\lambda,K)\le f^{\LOCC}(\lambda,K)
 \le f^{\cPPT}(\lambda,K)\le H_NT_K,
 \qquad H_N=\sum_{i=1}^N\frac1i.
\end{equation}
For $1<\alpha\le2$ one also has
\begin{equation}\label{eq:holder-converse}
 f^{\cPPT}(\lambda,K)
 \le\beta^{(2-\alpha)/\alpha}(N/K^2)^{(\alpha-1)/\alpha}
 \left(\sum_i\lambda_i^\alpha\right)^{1/\alpha}.
\end{equation}
At $\alpha=1$ the bound is $f\le\beta$.
\end{theorem}
\begin{proof}
The decreasing eigenvalues $\mu_r$ of a fidelity effect satisfy $\mu_r\le v_r$, by \eqref{eq:budget}. The spectral trace inequality and $\lambda_i\le L_i/i$ give
\[
 \Tr(MU\rho U^*)\le\sum_i\lambda_i\mu_i
 \le\sum_i\lambda_iv_i\le H_NT_K.
\]
For a prefix rank $r$ put $w=\min(a,K,\sqrt{N/r})$ and $d=\max(1,\lfloor w/2\rfloor)$. Then $d\ge w/4$ and \eqref{eq:packing-feasible} holds. Indeed $d=1$ is always feasible; otherwise $\lfloor a/d\rfloor\lfloor b/d\rfloor\ge N/(4d^2)\ge N/w^2\ge r$. Apply Lemma~\ref{lem:packing} to the span of the top $r$ eigenvectors. The remaining positive input contributes nonnegatively, giving fidelity at least $(d/K)L_r\ge v_rL_r/4$.

For \eqref{eq:holder-converse} let $q=\alpha/(\alpha-1)\ge2$. The effect eigenvalues satisfy $\Tr M^q\le\beta^{q-2}\Tr M^2$. Schatten Holder and \eqref{eq:budget} give the result.
\end{proof}

\section{One common protocol from polynomial orbit spaces}\label{sec:universal}

We give an intrinsic construction on the original tensor coordinates. It needs neither a supplied eigenbasis nor a Schur decomposition. All spaces, bases and integer packings below are fixed before the input state is quantified.

\subsection{Universal orbit projections}
Write $\mathcal W_{D,k}=\{1,\ldots,D\}^k$, and let $e_w$ denote a computational word. For $E\subseteq\mathcal W_{D,k}$ define
\begin{equation}\label{eq:orbit-space}
 V_k(E)=\operatorname{span}\{M^{\otimes k}e_w:M\in\C^{D\times D},\ w\in E\},
 \qquad B_{D,k}=\binom{k+D^2-1}{D^2-1}.
\end{equation}
Here the generating matrices are arbitrary complex matrices, not only unitaries. Let $P_k(E)$ be the orthogonal projection onto this known space.

\begin{lemma}[Polynomial orbit dimension and captured mass]\label{lem:orbit}
The projections are nested with $E$, commute with every $M^{\otimes k}$, and obey
\begin{equation}\label{eq:orbit-bound}
 \rank P_k(E)\le B_{D,k}|E|,\qquad
 \Tr[P_k(E)\rho^{\otimes k}]\ge\sum_{w\in E}p_{w_1}\cdots p_{w_k},
\end{equation}
where $p$ is any eigenvalue labeling of $\rho$. If $E\ne\varnothing$, then $V_k(E)$ contains $\operatorname{Sym}^k(\C^D)$.
\end{lemma}
\begin{proof}
For fixed $w$, the vector $M^{\otimes k}e_w$ is a homogeneous degree-$k$ polynomial in the $D^2$ entries of $M$. More precisely it is the image, under a fixed linear contraction, of $(\operatorname{vec}M)^{\otimes k}$. The span of these latter vectors has dimension $B_{D,k}$. Summing the resulting dimension bound over $E$ proves the rank estimate.

Composition of matrices makes $V_k(E)$ invariant under every $M^{\otimes k}$, including the adjoint powers. It is therefore a reducing subspace and its orthogonal projection commutes with those powers. If $U$ diagonalizes $\rho$, then $P_k(E)$ fixes $U^{\otimes k}e_w$ for $w\in E$. Its other diagonal entries in that eigenbasis are nonnegative, which proves the mass bound. Finally, for any $w\in E$ and vector $z$, choosing every column of $M$ equal to $z$ gives $M^{\otimes k}e_w=z^{\otimes k}$. Such powers span the symmetric subspace.
\end{proof}

Let $\widehat p_w$ be a word's empirical law and set
\[
 E_k(c)=\{w:H(\widehat p_w)\le c\},\qquad
 Q_{D,k}=B_{D,k}(k+1)^D.
\]
There are at most $(k+1)^D$ types, each containing at most $e^{kH(q)}$ words, so
\begin{equation}\label{eq:dimension-bounds}
 |E_k(c)|\le(k+1)^D e^{kc},\qquad
 \rank P_k(E_k(c))\le Q_{D,k}e^{kc}.
\end{equation}
These estimates hold for all real $c$, with the empty event allowed. For a law $p$, let $A_k(p,c)=\Prb_p\{H(\widehat p)\le c\}$.

\subsection{Geometric packing and a reserved pure sector}
Put
\begin{equation}\label{eq:reservation}
 \begin{aligned}
 J&=(k+1)^2,\qquad r_k=\binom{k+D-1}{D-1},\qquad Q=4Jr_k,\\
 a_0&=\lfloor a/Q\rfloor,\quad b_0=\lfloor b/Q\rfloor,
 \quad a_*=a-Ja_0,\quad b_*=b-Jb_0.
 \end{aligned}
\end{equation}
If $a_0=0$, use the unreserved version of the construction below: replace $a_0,b_0$ in \eqref{eq:grid} by $\lfloor a/J\rfloor,\lfloor b/J\rfloor$, assign the successive differences of $P_j$ directly to these $J$ sector pairs, and omit the separate symmetric sector. Extend the resulting isometry and complete the local decoder on unused coordinates as below. If either sector dimension is zero, use the identity preprocessing and complete local reset channels. This specifies a physical family at every finite block and every positive integer target. The following bounds concern $a_0>0$, which holds eventually for fixed dimensions.

For $0\le j<J$ define
\begin{equation}\label{eq:grid}
 d_j=\max\{1,\lfloor\min(K,a_0)/2^j\rfloor\},\qquad
 c_j=\frac{\log(a_0b_0/(4d_j^2))-\log Q_{D,k}}k.
\end{equation}
The $d_j$ decrease and the $c_j$ increase. Let $P_0=0$ and $P_{j+1}=P_k(E_k(c_j))$. By \eqref{eq:dimension-bounds},
\begin{equation}\label{eq:grid-budget}
 \rank P_{j+1}\le\frac{a_0b_0}{4d_j^2}
 \le\lfloor a_0/d_j\rfloor\lfloor b_0/d_j\rfloor.
\end{equation}
The last inequality follows from $1\le d_j\le a_0\le b_0$ and $\lfloor x\rfloor\ge x/2$ for $x\ge1$.

Let $H_k$ be the symmetric projector. Define $R_j=P_j-H_k$ when $P_j\ne0$, and $R_j=0$ otherwise. Lemma~\ref{lem:orbit} shows that these are nested projections orthogonal to $H_k$. Assign each orthogonal difference $R_{j+1}-R_j$ to its own pair of $a_0,b_0$ local sectors and use Bell width $d_j$. Its rank satisfies \eqref{eq:grid-budget}. Assign $H_k$ to the remaining $a_*,b_*$ sectors, with Bell width
\[
 d_* = \min\{K,d_{a_*,b_*}(r_k)\}.
\]
The reservation $Q=4Jr_k$ ensures $r_k\le a_*b_*$ and $d_{a_*,b_*}(r_k)\ge a_0$, so $d_*\ge d_0$. Indeed each large local sector contains at least $(4r_k-1)$ times its small local sector; width $a_0$ is consequently feasible for the symmetric rank $r_k$.

Apply Lemma~\ref{lem:packing} coherently to all assigned source sectors. Their images are orthogonal; their isometry extends to one unitary on the existing $N$-dimensional input. Sector-specific discard maps, together with reset Kraus maps on every remaining local coordinate, give two complete local channels. Orthogonal sector labels remove all cross-sector terms; commutation with $\rho^{\otimes k}$ also removes coherences with the unused source complement. This construction defines our single family $F_k^{\rm u}$.

\begin{proposition}[All prefixes and the pure promise in one family]\label{prop:prefix}
For every input density matrix and every $0\le j<J$,
\begin{equation}\label{eq:good-mass}
 F_k^{\rm u}(\rho,K)\ge\frac{d_j}K A_k(p,c_j).
\end{equation}
For every state $Y$ supported on $H_k$, the same protocol has fidelity at least
\[
 \min\{1,d_{a_*,b_*}(r_k)/K\}\Tr Y.
\]
\end{proposition}
\begin{proof}
Write $w_j=d_j/K$ and $B_j=\Tr[(H_k+R_j)\rho^{\otimes k}]$. The sector construction gives
\[
 F_k^{\rm u}\ge(d_*/K)B_0+\sum_{j=0}^{J-1}w_j(B_{j+1}-B_j).
\]
Since $d_*/K\ge w_0$, summation by parts makes its right side at least
\[
 w_{J-1}B_J+\sum_{j=1}^{J-1}(w_{j-1}-w_j)B_j.
\]
Every coefficient is nonnegative. Lemma~\ref{lem:orbit} bounds $B_{j+1}$ below by $A_k(p,c_j)$, and the corresponding classical cumulative masses increase. The resulting weighted sum is at least $w_jA_k(p,c_j)$ for every $j$. The pure-sector statement follows directly from coherent packing; it uses the same unitary and local channels.
\end{proof}

For fixed $m,n$, all reservation and orbit factors have logarithm $O_{m,n}(\log k)$. The removed local fraction is at most $J/Q=1/(4r_k)\to0$, so $a_*/a,b_*/b\to1$. An efficient implementation for arbitrary mixed inputs remains a separate question.

\subsection{An information-spectrum sandwich}
Sample $X$ with law $p$ on its positive support, put $Z=-\log p_X$, and let $S_k=Z_1+\cdots+Z_k$. Define the exact overhead
\begin{equation}\label{eq:overhead}
 h_{D,k}=\log(16Q^2Q_{D,k}).
\end{equation}

\begin{theorem}[Dimension-explicit finite-block sandwich]\label{thm:sandwich}
Whenever $K\le a_0$, put $x=\log(N/K^2)$. For every $t\ge0$,
\begin{equation}\label{eq:sandwich}
 \Prb\{S_k\le x-h_{D,k}\}\le F_k^{\rm u}(\rho,K)
 \le f_k^{\calC}(\rho,K)\le F_k^\star(\rho,K)
 \le\Prb\{S_k\le x+t\}+e^{-t/2},
\end{equation}
for each physical operation class containing product channels and contained in $\cPPT$. Zero eigenvalues are permitted. If $D\ge2$ and $k\ge1$, then
\begin{equation}\label{eq:overhead-explicit}
 0\le h_{D,k}\le10D^2\log(k+D).
\end{equation}
\end{theorem}
\begin{proof}
For $d_0=K$, Proposition~\ref{prop:prefix} gives $F_k^{\rm u}\ge A_k(p,c_0)$. Since $a_0\ge a/(2Q)$ and $b_0\ge b/(2Q)$, one has $kc_0\ge x-h_{D,k}$. For every positive-probability word,
\[
 kH(\widehat p_w)\le-\log p^{\otimes k}(w),
\]
by nonnegativity of relative entropy. This proves the lower bound.

For the converse, diagonalize the input and absorb the arbitrary preprocessing unitary into the fidelity effect. Its diagonal entries $z_i$ satisfy $0\le z_i\le1$ and $\sum_i z_i^2\le N/K^2$. With $r=N/K^2$, eigenvalues below $e^{-\log r-t}$ contribute at most
\[
 \sqrt{\sum_{\rm small}(p_i^{(k)})^2\sum_{\rm small}z_i^2}
 \le\sqrt{e^{-\log r-t}r}=e^{-t/2}.
\]
The remaining contribution is bounded by the event in \eqref{eq:sandwich}.

For \eqref{eq:overhead-explicit}, use $B_{D,k}\le(k+D^2)^{D^2-1}$, $r_k\le(k+D)^{D-1}$ and $\log(k+D^2)\le2\log(k+D)$. Substitution into \eqref{eq:overhead}, with $Q=4(k+1)^2r_k$, gives the displayed conservative constant, retaining every floor loss above.
\end{proof}

\begin{corollary}[Finite Gaussian error and shrinking variance]\label{cor:berry}
Let $V=\Var Z>0$, $T=\E|Z-H(p)|^3$, $\sigma=\sqrt{kV}$,
\[
 \beta=\frac{C_{\rm BE}T}{V^{3/2}\sqrt k},\qquad
 y=\frac{\log(N/K^2)-kH(p)}\sigma.
\]
There is an absolute Berry--Esseen constant $C_{\rm BE}$ such that, under the integer budget $K\le a_0$, both $F_k^{\rm u}$ and each $f_k^{\calC}$ in \eqref{eq:sandwich} differ from $\Phi(y)$ by at most
\begin{equation}\label{eq:gaussian-error}
 \beta+\frac{\max(h_{D,k},t)}{\sqrt{2\pi}\,\sigma}+e^{-t/2}
 \qquad(t\ge0).
\end{equation}
For changing dimensions and input laws, sufficient conditions are the eventual integer budget and
\begin{equation}\label{eq:triangular}
 kV\to\infty,\qquad T/(V^{3/2}\sqrt k)\to0,\qquad
 h_{D,k}/\sqrt{kV}\to0.
\end{equation}
In particular these follow if $DV$ is eventually bounded below by a positive constant, $DT$ is eventually bounded above, and $D^5\log^2(k+D)=o(k)$. The moving threshold $y$ need not converge.
\end{corollary}
\begin{proof}
Apply the finite iid Berry--Esseen inequality to both probability terms in \eqref{eq:sandwich}; the normal CDF has Lipschitz constant $1/\sqrt{2\pi}$. For \eqref{eq:triangular} choose $t=\sqrt\sigma$. This makes the smoothing ratio and exponential remainder vanish. For the concrete condition, $\sqrt{kV}\ge\sqrt c\sqrt{k/D}$ when $DV\ge c>0$, while $T/V^{3/2}\le C\sqrt D/c^{3/2}$ when $DT\le C$. Equation~\eqref{eq:overhead-explicit} reduces the overhead ratio to a constant times $D^{5/2}\log(k+D)/\sqrt k$. Its square tends to zero by hypothesis, and the same hypothesis forces $k/D\to\infty$. These are sufficient conditions for this construction, not necessary conditions or automatic APPT moment properties. Standard Berry--Esseen and source-coding results are credited in \cite{shevtsova,hayashi-second}.
\end{proof}

\section{Exact asymptotic extraction laws}\label{sec:asymptotic}

The limits in this section fix $m,n$ and an input state $\rho$. Let $p$ be its eigenvalue law, let $s=\rank\rho$, and put
\[
 A=\log m,\quad L=\log D,\quad
 C(p)=\min\left(A,\frac{L-H(p)}2\right),\quad
 C_0(p)=\min\left(A,\frac{L-\log s}2\right).
\]
Write $S_\alpha(p)=(1-\alpha)^{-1}\log\sum_i p_i^\alpha$, with $S_1=H$, and take $K_k=\lceil e^{kR}\rceil$ for a fixed $R\ge0$. At $R=0$ the target is one dimensional and fidelity is exactly one.

\begin{theorem}[One all-state family with the optimal laws]\label{thm:universal-laws}
The single family constructed in Section~\ref{sec:universal} has the following properties for every fixed input state, including rank-deficient states.
\begin{enumerate}
\item The faithful extraction capacities for $\LO$, $\LOCC$, and $\cPPT$ all equal $C(p)$. The common family achieves fidelity tending to one for each $R<C(p)$.
\item At every rate, its fidelity exponent equals that of the state-aware optimum in each of the three operation classes:
\begin{equation}\label{eq:exact-exponent}
 \lim_{k\to\infty}-\frac1k\log F_k^{\rm u}(\rho,K_k)
 =\lim_{k\to\infty}-\frac1k\log f_k^{\calC}(\rho,K_k)
 =E_p(R),
\end{equation}
where
\begin{equation}\label{eq:exponent-formula}
 E_p(R)=\pos{R-A}+\max_{1\le\alpha\le2}
 \frac{\alpha-1}{\alpha}\bigl(S_\alpha(p)-c_R\bigr),
 \qquad c_R=L-2\min(R,A).
\end{equation}
The summand at $\alpha=1$ is zero. In particular $R>C(p)$ gives a strictly positive strong-converse exponent.
\item In the strict direct regime $C_0(p)<R<C(p)$, put $c=L-2R$. Then
\begin{equation}\label{eq:reliability}
 \begin{aligned}
 \lim_{k\to\infty}-\frac1k\log(1-F_k^{\rm u})
 &=\lim_{k\to\infty}-\frac1k\log(1-F_k^\star)=I_p(c),\\
 I_p(c)&:=\min_{H(q)\ge c}D(q\|p)
 =\sup_{0<\alpha<1}\frac{1-\alpha}{\alpha}(c-S_\alpha(p)).
 \end{aligned}
\end{equation}
Relative entropy is evaluated on the positive support of $p$.
\item For $R<C_0(p)$ the same universal protocol has fidelity exactly one for all sufficiently large $k$. The eventual-exact capacity is $C_0(p)$, even for state-aware completely PPT postprocessing.
\end{enumerate}
\end{theorem}

\begin{proof}[Proof of the fidelity exponent]
The converse follows from \eqref{eq:holder-converse} on tensor powers and $k^{-1}\log K_k\to R$. For achievability assume first $R>0$ and fix a law $q$ on the positive support of $p$ with $H(q)<L$. Choose
\[
 0<b<\min\{R,A,(L-H(q))/2\},\qquad d_k=\lceil e^{kb}\rceil.
\]
For large $k$, $d_k\le\min(K_k,a_0)$, and the geometric grid contains a width between $d_k/2$ and $d_k$. Indeed the required grid index is $O_{m,n}(k)$, whereas $J=(k+1)^2$. Its entropy cutoff tends to $L-2b>H(q)$. An empirical type $q_k\to q$ is therefore included eventually; its probability is at least $(k+1)^{-D}e^{-kD(q_k\|p)}$. Proposition~\ref{prop:prefix} yields a fidelity exponent at most $D(q\|p)+R-b$. Let $b$ approach the displayed minimum and approximate the equality cases in entropy by laws on the same positive support. The only full-entropy boundary is the uniform full-support law, which is approached by nonuniform laws. Rate zero has target one and fidelity exactly one. Consequently
\begin{equation}\label{eq:variational-exponent}
 \limsup_k-\frac1k\log F_k^{\rm u}
 \le\min_q\left\{D(q\|p)+\max\left(0,R-A,R-L/2+H(q)/2\right)\right\}.
\end{equation}
All orbit, reservation and rounding factors have logarithm $O_{m,n}(\log k)$; the strict width approximation above removes them without a diagram or ordering assumption on $q$. The maximum in \eqref{eq:variational-exponent} is $\pos{R-A}+\pos{H(q)-c_R}/2$.

For completeness the matching scalar identity, with no exchange of minimax operations, is
\begin{equation}\label{eq:escort-identity}
 \min_q\{D(q\|p)+\pos{H(q)-c}/2\}
 =\max_{1\le\alpha\le2}\frac{\alpha-1}{\alpha}(S_\alpha(p)-c),\qquad c\ge0.
\end{equation}
Let $q_{\alpha,i}=p_i^\alpha/Z_\alpha$, $Z_\alpha=\sum_i p_i^\alpha$, and call the maximand $J_c(\alpha)$. Direct differentiation gives
\[
 J'_c(\alpha)=\frac{H(q_\alpha)-c}{\alpha^2},\qquad
 \frac{d}{d\alpha}H(q_\alpha)=-\alpha\Var_{q_\alpha}(\log p_i)\le0.
\]
If $H(p)\le c$, use $q=p$, and both sides vanish. If $H(q_2)\ge c$, use $q=q_2$, giving $(S_2(p)-c)/2$. Otherwise an escort with $1<\alpha<2$ has $H(q_\alpha)=c$, and its relative entropy is $J_c(\alpha)$. Boundary cases agree. Conversely, for every $q$ and $1\le\alpha\le2$,
\[
 D(q\|p)=\alpha^{-1}D(q\|q_\alpha)
 +\frac{\alpha-1}{\alpha}(S_\alpha(p)-H(q)),
\]
while $\pos{H-c}/2\ge(\alpha-1)(H-c)/\alpha$. This proves \eqref{eq:escort-identity}. Since $c_R\ge\log(n/m)\ge0$, the identity applies and proves \eqref{eq:exact-exponent} by squeezing all three operation classes between the common $\LO$ protocol and the $\cPPT$ converse.
\end{proof}

\begin{proof}[Proof of faithful and eventual-exact capacities]
For $0<R<C(p)$ one has $R<A$, so eventually $K_k\le a_0$. In \eqref{eq:sandwich},
$x-h=k(L-2R)+O(\log k)>k(H(p)+\delta)$ for a fixed $\delta>0$. The ordinary law of large numbers implies fidelity tending to one. Rate zero is trivial. Above $C(p)$ the exponent in \eqref{eq:exponent-formula} is positive: either $R>A$, or $J'_{L-2R}(1)>0$.

If $R<C_0(p)$, every positive-probability word has empirical entropy at most $\log s$, while the first grid cutoff tends to $L-2R>\log s$. Thus the entire classical mass in \eqref{eq:good-mass} is one eventually, giving exact fidelity. Conversely, exact fidelity for a rank-$s^k$ state forces its rotated support into the eigenvalue-one space of $M$. Hence $s^k\le\Tr M^2\le N/K^2$ and $K\le a$. Capacities here are suprema of rates realized by actual positive-integer target sequences and actual protocol sequences. These sets contain zero, since every trace-preserving channel to the one-dimensional target has fidelity exactly one. The upper bounds and every strict subcritical achievable rate therefore identify the suprema even when $C(p)$ or $C_0(p)$ is zero; no boundary attainability is inferred.
\end{proof}

\begin{proof}[Proof of strict-direct reliability]
The hypotheses imply $H(p)<c<\log s$. Proposition~\ref{prop:prefix} at $d_0=K_k$ bounds the actual error above by the probability that the empirical entropy exceeds $c-o(1)$. The number of types is polynomial, and each type has probability at most $e^{-kD(q\|p)}$. This gives the lower bound $I_p(c)$ on the error exponent of the same common protocol. The strict interval permits approximation of a boundary minimizer by laws of entropy strictly above the requested threshold.

To prove optimality, fix $0<t<1$. At most $N/(K^2t^2)$ eigenvalues of $M$ exceed $t$. If $L_j$ is the mass of the largest $j$ input eigenvalues, capped at one beyond their support, the spectral trace inequality gives
\begin{equation}\label{eq:tail-converse}
 1-F_k^\star\ge(1-t)
 \left(1-L_{\lfloor N/(K^2t^2)\rfloor}(p^{\otimes k})\right).
\end{equation}
For integers $M_k$ with $k^{-1}\log M_k\to c\in(H(p),\log s)$, the optimal-prefix tail has exponent $I_p(c)$. Here is the needed type argument. Include all types with entropy at most $c-\delta$; their total cardinality is at most $(k+1)^s e^{k(c-\delta)}\le M_k$ eventually. The excluded mass is at most $(k+1)^s e^{-kI_p(c-\delta)}$. Conversely, a type approaching any $q$ with $H(q)>c$ has exponentially more words than $M_k$, and at least fraction $1-M_k/r_q$ of its equiprobable words must be excluded. Its mass is at least $(k+1)^{-s}e^{-kD(q\|p)}$. Approximate an entropy-$c$ minimizer and let $\delta\downarrow0$. With $t=1/2$, \eqref{eq:tail-converse} gives the matching upper bound on every error exponent.

The requisite continuity and the last expression in \eqref{eq:reliability} follow from escorts with $0<\alpha<1$. Since $H(p)<c<\log s$, the positive eigenvalues are nonuniform and a unique escort has entropy $c$. The identity
\[
 D(q\|p)=\alpha^{-1}D(q\|q_\alpha)
 +\frac{1-\alpha}{\alpha}(H(q)-S_\alpha(p))
\]
proves the variational formula, equality at that escort, and continuity in the strict interval.
\end{proof}

\begin{remark}[The exact zero-error boundary]
Formula~\eqref{eq:reliability} excludes $R=C_0(p)$. For a rank-four state on $4\times4$, $K_k=2^k$ permits exact packing of its $4^k$-dimensional support into $\Phi_{2^k}$ and garbage. Thus the state-aware error can be exactly zero at $R=\log2=C_0(p)$ even for nonuniform positive eigenvalues. A finite reliability expression at $c=\log4$ would miss this integer packing effect.
\end{remark}

\begin{theorem}[Entropy-limited Gaussian threshold]\label{thm:gaussian}
Suppose $0<C(p)<A$ and
\[
 V(p)=\sum_{p_i>0}p_i(-\log p_i-H(p))^2>0.
\]
For any target sequence with $\log K_k=kC(p)+z\sqrt k+o(\sqrt k)$,
\begin{equation}\label{eq:gaussian-limit}
 F_k^{\rm u}(\rho,K_k)\longrightarrow\Phi(-2z/\sqrt{V(p)}),\qquad
 F_k^\star(\rho,K_k)\longrightarrow\Phi(-2z/\sqrt{V(p)}).
\end{equation}
For each operation class $\calC$ above and fixed $0<\varepsilon<1$, let $K_{k,\calC}^{\rm opt}(\varepsilon)$ be the largest positive integer target for which an actual state-aware protocol in that class attains fidelity at least $1-\varepsilon$. These maxima exist for all sufficiently large $k$ and obey
\[
 \log K_{k,\calC}^{\rm opt}(\varepsilon)=kC(p)
 +\tfrac12\sqrt{kV(p)}\,\Phi^{-1}(\varepsilon)+o(\sqrt k),
\]
and the common family attains the same expansion.
\end{theorem}
\begin{proof}
The strict local-dimension gap ensures $K_k\le a_0$ eventually. Here $h_{D,k}=O(\log k)$ and $\log(N/K_k^2)=kH(p)-2z\sqrt k+o(\sqrt k)$. Apply the ordinary central limit theorem to the two sides of \eqref{eq:sandwich}, for example with $t=k^{1/4}$; alternatively use Corollary~\ref{cor:berry}. The same squeeze applies to the genuine state-aware supremum in every intermediate physical class. For fixed error, define each optimum as the maximum of the set of positive integers admitting the requested actual fidelity. Bracket that entire set above and below a fixed Gaussian quantile: the converse excludes every larger integer, and rounded smaller targets provide actual witnesses. The feasible integer sets are eventually nonempty and bounded, hence their natural suprema are attained maxima. This establishes both expansions without assuming target monotonicity of the particular universal family.
\end{proof}

Theorem~\ref{thm:gaussian} retains both strict hypotheses. The local-dimension cap and zero varentropy are not included by continuity. The next section treats a cap-boundary phenomenon of a different, minimax kind.

\section{The sharp cost of an unknown support}\label{sec:support}

\subsection{A common support and an adversarial average}
Fix $1\le s\le D$. Define the intrinsic common support
\begin{equation}\label{eq:rank-support}
 S_{k,s}=\operatorname{span}\{z_1\otimes\cdots\otimes z_k:
 \dim\operatorname{span}(z_1,\ldots,z_k)\le s\},\qquad r_{k,s}=\dim S_{k,s}.
\end{equation}
Every $\rho^{\otimes k}$ with $\rank\rho\le s$ is supported here. Let $P_s$ be a fixed rank-$s$ orthogonal projection on $\C^D$, and form the literal Haar average
\begin{equation}\label{eq:haar-average}
 \Omega_{k,s}=\int_{U(D)}(UP_sU^*/s)^{\otimes k}\,dU.
\end{equation}
Haar measure constructs an adversarial ensemble and is never implemented as preprocessing. The range of $\Omega_{k,s}$ is exactly $S_{k,s}$: a vector has zero quadratic form against this average precisely when it is orthogonal to every rotated tensor-product range. The integrand is continuous and nonnegative and Haar measure has full support. Those ranges span \eqref{eq:rank-support}. This also proves minimality of the common support, without assuming a character or eigenvalue formula.

\begin{theorem}[Support dimension and reciprocal average purity]\label{thm:flat-average}
For fixed $D$ and $1\le s<D$, put
\begin{equation}\label{eq:rank-constant}
 g=s(D-s),\qquad A_{D,s}=\prod_{1\le i\le s<j\le D}(j-i),\qquad
 c_{D,s}=s^{-g}/A_{D,s}.
\end{equation}
Then
\begin{equation}\label{eq:support-asymptotic}
 r_{k,s}\sim c_{D,s}s^kk^g,\qquad
 \Tr\Omega_{k,s}^2\sim(c_{D,s}s^kk^g)^{-1},\qquad
 r_{k,s}\Tr\Omega_{k,s}^2\longrightarrow1.
\end{equation}
\end{theorem}
\begin{proof}
Let $C:\C^s\to\C^D$ be the coordinate isometry, so $P_s=CC^*$. For Haar $U$ set
\[
 G(U)=C^*UP_sU^*C,\qquad X(U)=\det G(U)=|\det(C^*UC)|^2,
 \qquad Y(U)=\Tr G(U)/s.
\]
The matrix $G$ is a positive contraction. Write $D=s+t$ and $M_q=\E X^q$. The standard unitary-minor moment calculation gives
\begin{equation}\label{eq:minor-moment}
 M_q=\prod_{i=0}^{s-1}\frac{\binom{q+i}{q}}{\binom{q+t+i}{q}},
 \qquad q^gM_q\longrightarrow
 \prod_{i=0}^{s-1}\frac{(t+i)!}{i!}=A_{D,s}.
\end{equation}
Here is a recursive derivation. After fixing the first Haar column, align its projections into the $s$-dimensional coordinate space and its orthogonal complement by block unitaries. In the determinant expansion the remaining coisometry has $s-1$ rows in dimension $D-1$; averaging over the stabilizer gives the same smaller minor moment. The squared projected column length has moment $(s)_q/(D)_q$, as follows by expanding coordinate monomials on the complex unit sphere. Thus $M_{s,t}(q)=\binom{q+s-1}{q}\binom{q+t+s-1}{q}^{-1}M_{s-1,t}(q)$, with $M_{0,t}(q)=1$. This proves the product; cancellation of its finite factorial ratios proves the limit. We use these standard Haar calculations as ingredients, not as novelty claims; see also the general polynomial integration framework in \cite{collins-sniady}.

Haar invariance and the tensor trace identity give
\begin{equation}\label{eq:two-moments}
 s^k\Tr\Omega_{k,s}^2=\E Y^k.
\end{equation}
If $\lambda_1,\ldots,\lambda_s\in[0,1]$ are the eigenvalues of $G$, arithmetic--geometric mean and $1-\prod_i\lambda_i\le\sum_i(1-\lambda_i)$ show
\[
 X\le Y^s,\qquad Y\le(s-1+X)/s.
\]
Consequently $\E Y^k\ge M_{\lceil k/s\rceil}$. To match this bound, fix an integer $N\ge1$, put $b=s(N+1)$ and $\delta=1-[2(N+1)^2]^{-1}$. For $\delta\le x\le1$,
\[
 ((s-1+x)/s)^b\le x^N.
\]
For $u=1-x$, the logarithm of the left side divided by the right side vanishes at $u=0$ and has derivative
\[
 \frac{-1+[(N+1)-N/s]u}{(1-u)(1-u/s)}\le0
 \quad\text{for }0\le u\le[2(N+1)^2]^{-1}.
\]
This proves the comparison. Splitting according to $X\ge\delta$ and setting $q=\lfloor k/b\rfloor$ therefore gives the exact bound
\[
 \E Y^k\le M_{Nq}+((s-1+\delta)/s)^{bq}.
\]
The exponential remainder is negligible after multiplication by $k^g$. Equation~\eqref{eq:minor-moment} bounds the resulting limsup by $A_{D,s}[s(N+1)/N]^g$. Let $N\to\infty$ and combine with the lower bound. This proves $s^kk^g\Tr\Omega_{k,s}^2\to s^gA_{D,s}=1/c_{D,s}$.

It remains to match the support rank; purity alone supplies only its lower bound. For a positive contraction $G$ and $k\le sq$, its $U(s)$ tensor twirl satisfies
\begin{equation}\label{eq:twirl-lower}
 \int_{U(s)}(VGV^*)^{\otimes k}\,dV\ \ge\ (\det G)^q I.
\end{equation}
To prove this, diagonalize $G$ and cyclically permute its $s$ eigenvalues. On each tensor word the product of the $s$ resulting diagonal entries is $(\det G)^k$. Arithmetic--geometric mean bounds their average below by $(\det G)^q$, since $0\le\det G\le1$ and $k\le sq$. Applying the positive Haar twirl makes every cyclic conjugate have the same average, proving \eqref{eq:twirl-lower}.

The compression $C^{\otimes k*}\Omega_{k,s}C^{\otimes k}$ is fixed by this $U(s)$ twirl. Integrating \eqref{eq:twirl-lower} for the actual $G(U)$ gives
\[
 C^{\otimes k*}\Omega_{k,s}C^{\otimes k}\ge s^{-k}M_q I.
\]
The averaged matrix belongs to the permutation-operator algebra by the tensor commutant identity, so it commutes with every rotated $(UP_sU^*)^{\otimes k}$. Haar invariance transfers the same lower bound to each such range. An eigenvector of $\Omega_{k,s}$ with eigenvalue below $s^{-k}M_q$ must therefore be orthogonal to all these ranges. Their span is $S_{k,s}$, so every positive eigenvalue is at least $s^{-k}M_q$. Taking traces and then $q=\lceil k/s\rceil$ proves
\[
 r_{k,s}\le s^k/M_{\lceil k/s\rceil}.
\]
On the other side, Cauchy--Schwarz gives $r_{k,s}\Tr\Omega_{k,s}^2\ge1$. The already proved purity limit and \eqref{eq:minor-moment} squeeze $r_{k,s}/(s^kk^g)$ to $c_{D,s}$. Multiplying the two limits proves the last assertion. No concentration hypothesis for Young diagrams or supplied spectral certificate is used.
\end{proof}

At $s=1$, these identities reduce to
\[
 S_{k,1}=\operatorname{Sym}^k(\C^D),\qquad
 r_{k,1}=\binom{k+D-1}{D-1},\qquad
 \Omega_{k,1}=\Pi_{\rm sym}/r_{k,1}.
\]
At $s=D$ one instead has exactly $r_{k,D}=N$ and $\Omega_{k,D}=I/N$.

\subsection{A target-uniform minimax law}
Define the finite, explicitly computable quantities
\begin{equation}\label{eq:dw}
 d_k=d_{a,b}(r_{k,s}),\qquad
 w_k=\min\left(a,\sqrt{N/r_{k,s}}\right),\qquad
 \eta_k=\sqrt{r_{k,s}\Tr\Omega_{k,s}^2}.
\end{equation}
Proposition~\ref{prop:average} gives for every finite $k$ and every integer $K\ge1$
\begin{equation}\label{eq:minimax-finite}
 \begin{aligned}
 \min(1,d_k/K)&\le\G_{\LO}(k,K;s)\le\G_{\LOCC}(k,K;s)
 \le\G_{\cPPT}(k,K;s),\\
 \G_{\cPPT}(k,K;s)&\le\min\left(1,\frac aK,\frac{\sqrt{N\Tr\Omega_{k,s}^2}}K\right).
 \end{aligned}
\end{equation}
The lower construction knows $s$, but no property of the particular input within the rank-budget class.

\begin{theorem}[Entire sharp minimax profile]\label{thm:minimax}
For every fixed $2\le m\le n$ and $1\le s<D$, and every $\calC\in\{\LO,\LOCC,\cPPT\}$,
\begin{equation}\label{eq:minimax-uniform}
 \sup_{K\in\N,\ K\ge1}
 \left|\frac{\G_{\calC}(k,K;s)}{\min(1,w_k/K)}-1\right|\longrightarrow0.
\end{equation}
An explicit finite-$k$ bound for the supremum is
\begin{equation}\label{eq:minimax-relative-bound}
 \max\{1-d_k/w_k,\eta_k-1\}.
\end{equation}
Thus the relative equivalence includes every target integer, even when fidelity is exponentially small.
\end{theorem}
\begin{proof}
It remains to prove $d_k/w_k\to1$, since $\eta_k\to1$ by Theorem~\ref{thm:flat-average}. If $n>ms$, then
$r_{k,s}/(b/a)\to0$ and eventually $d=a$ is feasible. Consequently $d_k=w_k=a$ eventually.

If $n\le ms$, Theorem~\ref{thm:flat-average} gives eventually
\[
 w_k=\sqrt{N/r_{k,s}},\qquad
 w_k/a\to0,\quad w_k/b\to0,\quad w_k\to\infty.
\]
The polynomial factor is essential when $n=ms$. Put $q_k=w_k/a+w_k/b$, $\epsilon_k=\sqrt{q_k}$, and $d=\lfloor w_k/(1+\epsilon_k)\rfloor$. For large $k$ this is a positive integer with $d/w_k\to1$. Dividing a lower bound for the product of floors by $r=N/w_k^2$ gives
\[
 \frac{\lfloor a/d\rfloor\lfloor b/d\rfloor}{r}
 \ge (w_k/d)^2-(w_k/d)q_k+1/r.
\]
For $0<q_k<1$ the function $z^2-zq_k$ is increasing on $z\ge1$, and $w_k/d\ge1+\sqrt{q_k}$. The last display is therefore at least
$1+2\sqrt{q_k}-q_k^{3/2}>1$. Thus $d$ is feasible. Also every feasible integer is at most $w_k$, so $d_k/w_k\to1$.

Finally \eqref{eq:minimax-finite} bounds the ratio in \eqref{eq:minimax-uniform} between $d_k/w_k$ and $\eta_k$, independently of $K$. This proves both assertions.
\end{proof}

For $s=D$ the exact minimax value is $1/K$ at every blocklength: the maximally mixed state gives this upper bound and a product output attains it. Consequently its largest exact target is one, and, for $0\le\varepsilon<1$, its largest target with worst-class fidelity at least $1-\varepsilon$ is exactly $\lfloor(1-\varepsilon)^{-1}\rfloor$.

\begin{corollary}[Exact and fixed-error target sizes]\label{cor:targets}
Let $K_{\calC}^0(k;s)$ be the largest target dimension admitting exact fidelity for every state in the rank-budget class. Then
\begin{equation}\label{eq:exact-target-bound}
 d_k\le K_{\LO}^0\le K_{\LOCC}^0\le K_{\cPPT}^0\le\lfloor w_k\rfloor.
\end{equation}
If $n\le ms$ and $s<D$,
\begin{equation}\label{eq:logarithmic-cost}
 \log K_{\calC}^0(k;s)
 =\frac k2\log(D/s)-\frac{s(D-s)}2\log k-\frac12\log c_{D,s}+o(1).
\end{equation}
If $n>ms$, instead $K_{\calC}^0(k;s)=m^k$ eventually. For $s<D$ and fixed $0<\varepsilon<1$, the largest target admitting worst-class fidelity at least $1-\varepsilon$ satisfies
\begin{equation}\label{eq:fixed-error-minimax}
 \log K_{\calC}^{\varepsilon}(k;s)
 =\log w_k-\log(1-\varepsilon)+o(1).
\end{equation}
\end{corollary}
\begin{proof}
Exact performance gives $\Tr[\Omega_{k,s}(I-U^*MU)]=0$. Positivity and strict positivity of $\Omega_{k,s}$ on $S_{k,s}$ force $U^*MU$ to be the identity on that subspace. Hence $r_{k,s}\le\Tr M^2\le N/K^2$ and $K\le a$. This proves the upper bound in \eqref{eq:exact-target-bound}; packing proves its lower bound. Use $d_k/w_k\to1$ and \eqref{eq:support-asymptotic} for \eqref{eq:logarithmic-cost}. For fixed error, the actual lower construction works at $K=\lfloor d_k/(1-\varepsilon)\rfloor$, while \eqref{eq:minimax-finite} bounds every feasible target by $\eta_kw_k/(1-\varepsilon)$. Since $w_k\to\infty$ for $s<D$, integer rounding contributes $o(1)$ to the logarithm.
\end{proof}

The boundary $n=ms$ has first-order capacity $\log m$, but still incurs the logarithmic correction in \eqref{eq:logarithmic-cost}. The coefficient $s(D-s)/2$ is one quarter of the real dimension of the Grassmannian of $s$-dimensional supports. This interpretation follows from the two-moment proof; parameter counting is not used as a substitute for it.

\begin{corollary}[A known spectrum with unknown orientation]\label{cor:known-spectrum}
Fix any positive rank-$s$ spectrum. Requiring one exact protocol for its entire global-unitary orbit gives the same bounds \eqref{eq:exact-target-bound} and hence the same exact-target asymptotic. If the known spectrum is flat on rank $s$, the full target-uniform minimax profile \eqref{eq:minimax-uniform} also holds for that orbit alone.
\end{corollary}
\begin{proof}
For any positive rank-$s$ spectrum, each rotated tensor state is strictly positive on the tensor power of its rotated $s$-dimensional support. Those ranges span exactly $S_{k,s}$, so its orbit average has that range. The exact converse and common-subspace packing are unchanged. For the flat spectrum its average is precisely \eqref{eq:haar-average}, so the entire previous minimax argument applies. No fixed-error profile for an arbitrary nonuniform spectrum is inferred from this support statement.
\end{proof}

\subsection{Balanced unknown pure inputs and an exact circuit}
When $m=n$, $D=m^2$, and $s=1$, put $r_k=\binom{k+D-1}{D-1}$. Theorem~\ref{thm:minimax} gives, uniformly in every target integer,
\[
 \G_{\calC}(k,K;1)\sim\min\left(1,\frac{m^k}{K\sqrt{r_k}}\right).
\]
In particular
\begin{equation}\label{eq:pure-full-output}
 \G_{\calC}(k,m^k;1)\sim r_k^{-1/2}
 \sim\sqrt{(D-1)!}\,k^{-(D-1)/2}.
\end{equation}
A known pure state admits fidelity one at this same target. The common-protocol loss is therefore sharp and invisible to the first-order capacity.

The symmetric compression mechanism is established prior work \cite{plesch,baertschi,chen}; coherent streaming concentration also predates this work \cite{streaming}. The following specialization records its exact resource conversion in our model.

\begin{proposition}[Optimal integer Bell pairs with no supplied memory]\label{prop:pure-circuit}
For $m=n=2$ and every $k\ge1$, one circuit on the original $2k$ qubits, followed by local discarding, extracts exactly
\[
 h_k=\left\lfloor k-\frac12\log_2\binom{k+3}{3}\right\rfloor
\]
Bell pairs from every unknown pure iid input. This integer is optimal even for a common global unitary followed by completely PPT postprocessing. The circuit can be synthesized with $O(k^7\log^2(k+3))$ CNOT and arbitrary single-qubit gates, with no qubit supplied, reset, measured, or discarded before the global circuit. For any fixed finite inverse-closed library whose generated subgroup is dense in $SU(2)$, and every $\delta>0$, there is a circuit using that library and CNOT on those same wires whose operator-norm error, up to global phase, is at most $\delta$, whose length is at most $C G\,[\log_2\max(1,G/\delta)+2]^5$, and whose output fidelity is at least $1-2\delta$. Here $G$ is the larger of one and the actual continuous-gate circuit size, which has the preceding asymptotic bound, and $C$ depends only on the fixed library. Proposition~\ref{prop:effective-compiler} supplies an actual compiler under numerical access to the fixed library. A classical finite-precision synthesis-time bound remains separate from the circuit-size guarantee.
\end{proposition}
\begin{proof}
Let $r_j=\binom{j+3}{3}$ and $\ell_j=\lceil\log_4 r_j\rceil$. Fix an enumeration $\operatorname{enc}_j$ of the four-component occupations of $j$. Sequentially compress the first $j$ four-level systems into $\ell_j$ of those same systems; all other processed systems become zero on the entire symmetric subspace. This is possible because $r_j/r_{j-1}=(j+3)/j\le4$, so $\ell_j\le\ell_{j-1}+1\le j$.

At step $j$, the only active systems are the $\ell_{j-1}$ memory systems and the next original input. In their computational basis the symmetric branching vectors are
\[
 v_\nu=\sum_{i:\nu_i>0}\sqrt{\nu_i/j}
 |\operatorname{enc}_{j-1}(\nu-e_i)\rangle|i\rangle.
\]
Their computational supports are disjoint, and each has at most four entries. At most three real Givens rotations per occupation send it to one selected basis state. Complete these states to the desired $\operatorname{enc}_j(\nu)$ labels by a fixed permutation of the full active basis. If the active register has $q_j=2(\ell_{j-1}+1)$ qubits and dimension $L_j=2^{q_j}$, at most $4L_j$ two-level operations suffice, and $L_j\le16r_{j-1}=O(j^3)$. All rotations and the permutation are computable from polynomially many occupation labels, not from the input state.

For clarity an elementary ancilla-free synthesis bound for a two-level operation on $q$ qubits is $O(q^2 2^q)$. A Gray path reduces it to pattern-controlled one-qubit rotations and phases. A phase on one selected string is $\exp(i\theta P_x)$ with
\[
 P_x=2^{-q}\prod_l(I+(-1)^{x_l}Z_l).
\]
Expand into at most $2^q$ commuting Pauli strings; implement each nonempty string by a parity circuit, a one-qubit rotation, and the inverse parity circuit, using an existing participating wire. Controlled $Y$ rotations have the same expansion after a local change of basis. A Gray path uses at most $2q$ selective transpositions. Thus a conservative total bound is
\[
 1000\sum_{j=2}^k q_j^2L_j^2+100k^2=O(k^7\log^2(k+3)).
\]
This agrees with the general synthesis context \cite{barenco,shende}, while making the absence of a supplied workspace explicit.

After compression at $j=k$, at most $2\ell_k=2(k-h_k)$ memory qubits can be nonzero. Redistribute them among the two garbage registers of $k-h_k$ qubits each. All remaining original qubits are already zero. Hadamards and cross-party CNOTs on the remaining $h_k$ pairs create the desired Bell states; local discarding removes the garbage. The zero wires were generated from the promised input by the unitary compression, not supplied externally.

For necessity, exact fidelity for every pure input makes the effect the identity on the whole symmetric subspace. Lemma~\ref{lem:rains} then gives $r_k\le4^k/4^h$, proving the exact integer optimum.

For the finite-library size statement, compactness and the density assumption give a fixed finite basic net. Every $SU(2)$ residual admits an exact balanced commutator $[A,B]$, with each factor at distance at most the square root of the residual's distance to the identity. Cancellation in an approximate commutator gives the error bound $4\delta\varepsilon+2\varepsilon^2$ when the exact factors are $\delta$-close to identity and their approximations have error $\varepsilon$. Recursive correction gives actual finite words of error $2^{-p}$ and length at most $C(p+1)^5$. These are the standard Solovay--Kitaev ingredients \cite{dawson}; the stated conservative power is sufficient here. Remove the overall phase of each original single-qubit gate, approximate its $SU(2)$ representative to error $\delta/G$, and place its word on the same original wire. The wire embedding preserves the operator norm. Telescoping over the whole original gate list gives error at most $\delta$ up to a total phase. The same actual local discard then has output fidelity at least $1-2\delta$.

\begin{proposition}[Effective original-register compiler]\label{prop:effective-compiler}
Fix a finite inverse-closed universal library $\{V_i:i\in[M]\}\subset SU(2)$, its inverse-label table, and numerical entry routines giving rational quaternion coordinates to every requested absolute precision $2^{-r}$. Fix a finite $1/576$-net of library words, each of length at most $L$; these are library-dependent offline data. There is a total compiler, uniform in the integers $k,p\ge0$, whose output is a finite list $W_{k,p}$ of library labels with original-wire indices and CNOT instructions on the original $2k$ wires. Let $N_k$ be the actual length of its exact continuous-gate descriptor list and set
\[
 q=p+1+\left\lceil\log_2\max(1,N_k)\right\rceil.
\]
The exact list has the polynomial size bound in the preceding theorem. The compiled list satisfies
\[
 |W_{k,p}|\le32(L+1)N_k(q+1)^5,
 \qquad \|U_k-\zeta_k\,\mathsf U(W_{k,p})\|\le2^{-(p+1)},
\]
Here $\mathsf U(W)$ denotes the unitary of an instruction list, $|\zeta_k|=1$, and $U_k$ is the actual exact extraction circuit. The same local discard therefore produces target overlap at least $1-2^{-p}$ for every unknown pure iid input. The numerical routines are the only oracle inputs; when they are computable, the compiler is computable. No bound on its classical bit running time is asserted.
\end{proposition}

\begin{proof}
Enumerate the occupation labels by bounded triples and use their finite-list rank and unrank maps for the memory code. At every stage, choose the first occupied pivot, order the remaining branches by their integer labels, and align the output code by explicit swaps. This produces the exact gate descriptor list using finite arithmetic and lists; no arbitrary complement enumeration is needed. Each descriptor records an integer occupation-mass ratio, a sign and a number of half-angle operations. Its sine and cosine have total rational approximation procedures, including the zero-mass cases, obtained from rational arithmetic, integer square root and half-angle identities. The one-qubit descriptor has a concrete $SU(2)$ representative times a unit-modulus phase. Multiplying these phases gives the actual $\zeta_k$.

Compactness and universality provide the fixed finite net. For each representative, use the numerical balanced-commutator recursion with two correction steps and overlapping rational stopping tests. Its finite cached program uses rational comparisons and conditioned chart selection; the normalized factor routines are computed internally. The proved numerical recursion returns a library word of error at most $2^{-q}$ and length at most $32L(q+1)^5$. Substitute these actual words for the one-qubit descriptors and retain each original CNOT. Telescoping over the actual $N_k$ descriptors gives error at most $N_k2^{-q}\le2^{-(p+1)}$. The factor $L+1$ retains the CNOTs in the length bound, and the definition using $\max(1,N_k)$ also covers an empty list. The output-effect perturbation bound for the same local discard is $2\|U_k-\zeta_k\mathsf U(W_{k,p})\|$, yielding the stated overlap uniformly in the input.
\end{proof}

The accompanying development also gives finite cached input-demand and arithmetic-expression operation counts, an actual common-denominator integer evaluator, and intermediate integer-size budgets. These statements distinguish computability and arithmetic cost from classical bit running time. The latter additionally requires a polynomial bound on every recursively generated operand and the costs of the actual numerical providers and structural operations.
\end{proof}

This proposition implements the pure-promise task. It does not provide an efficient implementation of the complete all-state mixed-input family of Section~\ref{sec:universal}.

\subsection{Compatibility with an all-state family}
The family in Section~\ref{sec:universal} already reserves the symmetric sector, so it retains the sharp balanced pure profile without a purity flag. Let $a=b=m^k$. In \eqref{eq:reservation} the remaining local dimensions satisfy $a_*/a,b_*/a\to1$. Integer packing of their symmetric rank $r_k$ then gives
\[
 \frac{d_{a_*,a_*}(r_k)}{a/\sqrt{r_k}}\longrightarrow1.
\]
For example $\lfloor a_*/\lceil\sqrt{r_k}\rceil\rfloor$ is feasible; $r_k\to\infty$ and $a/\sqrt{r_k}\to\infty$ remove its relative rounding loss. Proposition~\ref{prop:prefix} therefore proves, for every $\varepsilon>0$, one blocklength threshold such that simultaneously for all positive integer $K$ and every unknown pure input,
\[
 F_k^{\rm u}(\rho,K)\ge(1-\varepsilon)
 \min\left(1,\frac{m^k}{K\sqrt{r_k}}\right).
\]
The averaged-state converse supplies the matching upper profile for any uniform guarantee over pure inputs. The same selected family retains all the mixed-state laws, the finite Gaussian estimates and the fixed-error inverse above. The general rank-$s$ sharp minimax theorem retains its rank-budget promise; simultaneous sharp constants for every $s>1$ without such a promise are not inferred.

\section{Applications, verification, and open questions}\label{sec:applications}

\subsection{APPT entropy as an extraction budget}
A state on $\C^m\otimes\C^n$ is absolutely PPT if every global unitary conjugate has positive partial transpose. The capacity formula connects its spectral entropy to a collective extraction resource, even though no single-copy global rotation can produce an NPT state. Qualitative activation from purity is known \cite{kondra}; the following identity records the quantitative link needed here.

Let
\[
 E_1(m,n)=\log D-\min_{\rho\in\APPT(m,n)}S(\rho).
\]
The APPT set is compact: it is the intersection, inside the compact set of density matrices, of closed partial-transpose positivity conditions. Von Neumann entropy is continuous, so the minimum is attained. Theorem~\ref{thm:universal-laws} immediately gives
\begin{equation}\label{eq:appt-capacity}
 \max_{\rho\in\APPT(m,n)}C(p_\rho)
 =\min\left(\log m,\frac{E_1(m,n)}2\right).
\end{equation}
Thus a minimum-entropy theorem for APPT spectra calibrates the maximal capacity in the present model. A companion all-dimension entropy profile can be substituted in this identity \cite{zhang-appt-entropy}; its proof and its verification status are separate from the present paper.

One can obtain a self-contained operational comparison from the already verified purity asymptotic \cite{zhang-appt-purity}. For $P_{\max}(m,n)=\max_{\APPT}\Tr\rho^2$, Jensen's inequality gives $S(\rho)\ge-\log\Tr\rho^2$, and hence
\begin{equation}\label{eq:appt-purity-upper}
 \max_{\APPT}C\le\frac12\log(DP_{\max}).
\end{equation}
For balanced systems $D=m^2\to\infty$, the purity law $P_{\max}=D^{-1}+8D^{-2}+o(D^{-2})$ yields $\max_{\APPT}C\le4/D+o(D^{-1})$.

For a lower example take $\rho=(I+2vv^*)/(D+2)$. It is APPT for every unit vector $v$: after every global rotation, the smallest eigenvalue of $(vv^*)^\Gamma$ is at least $-1/2$, by the Schmidt coefficient inequality $t_it_j\le1/2$. Its spectrum is $(3,1,\ldots,1)/(D+2)$, and the local cap is inactive for every $2\le m\le n$. Hence
\[
 C(p_\rho)=\frac12\left[\log\frac D{D+2}+\frac{3\log3}{D+2}\right]
 =\frac{3\log3-2}{2D}+O(D^{-2}).
\]
With $c_{\rm sp}=(3\log3-2)/2$, an explicit error bound is
\[
 0\le c_{\rm sp}/D-C(p_\rho)\le\frac{3\log3}{D(D+2)}.
\]
Together with \eqref{eq:appt-purity-upper}, this gives, along every sequence of balanced dimensions tending to infinity,
\[
 c_{\rm sp}\le\liminf D\max_{\APPT}C
 \le\limsup D\max_{\APPT}C\le4.
\]
The actual maximum and the displayed comparison are established without solving the entropy optimizer. These are applications of the two-stage spectral capacity and do not assert ordinary fixed-state LOCC distillation at these rates.

Equation~\eqref{eq:appt-capacity} fixes the dimensions before taking the many-copy capacity limit. A joint growing-dimension conclusion requires uniform finite-block control. Theorem~\ref{thm:sandwich} supplies such control without spectral assumptions; a Gaussian conclusion additionally requires the moment and output-size conditions in \eqref{eq:triangular}. No assertion that every APPT near optimizer satisfies those conditions is made.

\subsection{Formalization and independent verification}
The development uses Lean~4.34.1 and a fixed Mathlib commit. Its objects are the original complex density matrices, arbitrary global unitaries, actual finite Kraus channels and partial transpose on the stated first party. Complete positivity is defined by positivity of every finite auxiliary amplification; the Choi criterion and actual complete Kraus representations are derived. Finite LOCC is an adaptive local-instrument tree with public outcomes and hidden Kraus indices. Closure LOCC is its genuine finite-dimensional coefficient closure, and its completely PPT inclusion follows from closedness of the Choi positivity and trace identities.

The coherent packing theorem includes unitary extension on the existing input, all unused-coordinate decoder branches and arbitrary logical coherences. The Rains converse is realized by the actual measure-and-prepare channel in Lemma~\ref{lem:rains}. Asymptotic endpoints use the selected common family itself and genuine physical suprema. Capacities and integer target optima are defined by actual protocol witnesses and success conditions; exact feasibility is not inferred from an unattained fidelity supremum. The probability development includes zero eigenvalues, finite Berry--Esseen control and true varying-law limits.

The revision uses intrinsic orbit projections in Section~\ref{sec:universal} and Haar minor moments plus an operator lower bound in Theorem~\ref{thm:flat-average}. Its finite overhead is \eqref{eq:overhead}, rather than the different Schur-block overhead in the preserved predecessor. The tensor commutant result is proved internally; a caller-supplied Young decomposition, moment identity, asymptotic rank law or optimal-performance certificate is not a hypothesis of these endpoints.

The verification archive separates the frozen mathematical completion from the effective compiler extension:
\begingroup
\renewcommand{\baselinestretch}{1.0}\small
\begin{center}
\begin{tabular}{>{\raggedright\arraybackslash}p{0.57\textwidth}>{\raggedright\arraybackslash}p{0.31\textwidth}}
\toprule
Mathematical component & Exact source scope\\
\midrule
Actual coherent packing and channel-to-effect bounds & Preserved scoped Lean audits\\
Complete positivity, Kraus equivalence, LOCC and Rains realization & Completion audit\\
One common family, capacities, exponents, Gaussian and integer targets & Completion audit\\
Actual Haar moments, rank, purity and physical minimax profiles & Completion audit\\
Original-register circuits and finite-library word-length bounds & Completion audit\\
APPT capacity extrema and explicit first-order comparisons & Completion audit\\
Computable descriptors, numerical compiler and actual output fidelity & Separate algorithm scope\\
Uniform classical compiler bit running time & Open; not a circuit-size claim\\
\bottomrule
\end{tabular}
\end{center}
\endgroup
The fixed-source completion audit contains 518 named theorem endpoints and 175 concrete model definitions. Its independent fresh build and empty-environment kernel replay cover a 72,820-declaration proof-and-definition closure at trust level zero. All 84 same-type forged-proof controls are rejected at their precise theorem roots after positive replay of their type dependencies. The only admitted foundational axioms are \texttt{propext}, \texttt{Classical.choice} and \texttt{Quot.sound}. A separate algorithm scope binds the new numerical program, computable descriptors, actual finite-label output and their concrete semantics. The exact declaration lists, source hashes, execution inputs, replay logs, control results and CI artifact checks are supplied as machine-readable evidence; source-map counts are not counts of discoveries.

The archived manifests identify the source to which each result applies. The older 608-root and APPT audits retain their original inputs. This manuscript's predecessor and the first working revision are preserved, and new results do not inherit their verification claims merely by being in the same repository. The source maps connect the named mathematical statements and physical models to the formal declarations. Formal verification certifies those precise types under their listed hypotheses; the correspondence with prose is independently reviewed rather than automatically inferred from a passing build. Internal AI review is not external human peer review or a certification of novelty.

\subsection{Implementation and further questions}
The pure-promise result supplies an actual computable finite-label compiler on the original register, with the whole-circuit error and output-fidelity guarantees of Proposition~\ref{prop:effective-compiler}. Its fixed net and inverse-label table depend on the library; no uniform procedure for finding a net from an arbitrary real-valued library description is claimed. A full classical bit-time theorem still requires complete operand-size, provider and structural-cost bounds. The general mixed-state construction remains an existence theorem for a specified finite unitary and local channels. Efficient representation transforms \cite{bch,kirby,krovi} provide related tools; a complete implementation must also account for reversible sector indexing, variable Bell dimensions and unitary completion on the available register.

The sharp minimax results fix $D$ and $s$. Their finite inequalities remain valid when these parameters vary, but the inverse-moment asymptotic requires additional uniform estimates before such a joint limit can be claimed. The explicit conditions in \eqref{eq:triangular} concern a different, information-spectrum route and do not replace those estimates. The exact reliability law at the zero-error boundary, cap-boundary Gaussian laws for arbitrary mixed inputs, and optimal circuit costs likewise remain open here.

\paragraph{Acknowledgments and disclosure.}
The author received no external funding. OpenAI Codex AI assistants were used throughout all six numbered sections and the accompanying Lean development for proof exploration, drafting and revision of mathematical arguments and prose, generation and debugging of formal code, and internal critical review. Responsibility for the mathematical statements and the distinction between written and machine-verified proofs remains with the author. Internal AI review is not external human peer review. Original materials retain all rights reserved; third-party results and software retain their own attribution and licenses.

\paragraph{Predecessors and reproducibility.}
This manuscript consolidates the author's repository preprints on state-aware and state-independent extraction \cite{zhang-extraction,zhang-universal}. Their original versions are preserved. The finite-block construction at all blocklengths, target-uniform rank-budget minimax theorem, support/average two-moment comparison, and exact pure-promise resource accounting are identified separately in the contribution and source maps. The reproduction bundle contains this manuscript, its sources, the mathematical verification records, and the separately scoped Lean packages; it does not turn earlier preprints into distinct additional discoveries.

\begin{thebibliography}{99}
\bibitem{rains} E.~M. Rains, ``A semidefinite program for distillable entanglement,'' \emph{IEEE Trans. Inf. Theory}, vol.~47, pp.~2921--2933, 2001. \href{https://arxiv.org/abs/quant-ph/0008047}{arXiv:quant-ph/0008047}.
\bibitem{lami} L. Lami, ``On PPT entanglement distillation,'' 2026. \href{https://arxiv.org/abs/2610.12454}{arXiv:2610.12454}.
\bibitem{jhhh} R. Jozsa, M. Horodecki, P. Horodecki, and R. Horodecki, ``Universal quantum information compression,'' \emph{Phys. Rev. Lett.}, vol.~81, p.~1714, 1998. \href{https://arxiv.org/abs/quant-ph/9805017}{arXiv:quant-ph/9805017}.
\bibitem{kw} M. Keyl and R.~F. Werner, ``Estimating the spectrum of a density operator,'' \emph{Phys. Rev. A}, vol.~64, p.~052311, 2001. \href{https://arxiv.org/abs/quant-ph/0102027}{arXiv:quant-ph/0102027}.
\bibitem{collins-sniady} B. Collins and P. \'{S}niady, ``Integration with respect to the Haar measure on unitary, orthogonal and symplectic group,'' \emph{Commun. Math. Phys.}, vol.~264, pp.~773--795, 2006. \href{https://arxiv.org/abs/math-ph/0402073}{arXiv:math-ph/0402073}.
\bibitem{hm} K. Matsumoto and M. Hayashi, ``Universal distortion-free entanglement concentration,'' \emph{Phys. Rev. A}, vol.~75, p.~062338, 2007. See also M. Hayashi and K. Matsumoto, \href{https://arxiv.org/abs/quant-ph/0109028}{arXiv:quant-ph/0109028}.
\bibitem{hayashi-second} M. Hayashi, ``Second-order asymptotics in fixed-length source coding and intrinsic randomness,'' \emph{IEEE Trans. Inf. Theory}, vol.~54, pp.~4619--4637, 2008. \href{https://arxiv.org/abs/cs/0503089}{arXiv:cs/0503089}.
\bibitem{hayashi-univ} M. Hayashi, ``Universal approximation of multi-copy states and universal quantum lossless data compression,'' \emph{Commun. Math. Phys.}, vol.~293, pp.~171--183, 2010. \href{https://arxiv.org/abs/0806.1091}{arXiv:0806.1091}.
\bibitem{takagi} R. Takagi, K. Watanabe, T. Matsuura, H. Arai, and M. Hayashi, ``Universal distillation of quantum entanglement,'' 2026. \href{https://arxiv.org/abs/2609.40215}{arXiv:2609.40215}.
\bibitem{tirone} S. Tirone, F.~A. Mele, V. Giovannetti, and L. Lami, ``Universal entanglement distillation,'' 2026. \href{https://arxiv.org/abs/2609.40080}{arXiv:2609.40080}.
\bibitem{llf} Z. Lin, K. Li, and K. Fang, ``Finite-blocklength and exponential analyses for universal entanglement and coherence distillation under various free operations,'' 2026, version~3. \href{https://arxiv.org/abs/2601.10190v3}{arXiv:2601.10190v3}.
\bibitem{rizzo} J. Rizzo, L. Lami, J. Eisert, and L. Leone, ``Universal quantum coding,'' 2026. \href{https://arxiv.org/abs/2609.38038}{arXiv:2609.38038}.
\bibitem{hayden-qmdl} P. Hayden, A. Maloney, J. Wang, and Y. Yang, ``Free entropy and quantum minimum description length,'' 2026. \href{https://arxiv.org/abs/2610.06741}{arXiv:2610.06741}.
\bibitem{watanabe} K. Watanabe and R. Takagi, ``Universal work extraction in quantum thermodynamics,'' \emph{Nat. Commun.}, vol.~17, art.~1857, 2026. \href{https://doi.org/10.1038/s41467-026-69143-3}{doi:10.1038/s41467-026-69143-3}.
\bibitem{zeng} S. Zeng, ``Reliability is not free in universal quantum work extraction,'' 2026. \href{https://arxiv.org/abs/2607.16702}{arXiv:2607.16702}.
\bibitem{plesch} M. Plesch and V. Bu\v{z}ek, ``Efficient compression of quantum information,'' \emph{Phys. Rev. A}, vol.~81, p.~032317, 2010. \href{https://arxiv.org/abs/0907.1764}{arXiv:0907.1764}.
\bibitem{baertschi} A. B\"artschi and S. Eidenbenz, ``Deterministic preparation of Dicke states,'' 2019. \href{https://arxiv.org/abs/1904.07358}{arXiv:1904.07358}.
\bibitem{chen} Z. Chen, C. Xiong, K. Khadiev, and J. Wu, ``Efficient quantum compression for identically prepared states with arbitrary dimension,'' version~2, 2026. \href{https://arxiv.org/abs/2307.12024v2}{arXiv:2307.12024v2}.
\bibitem{streaming} R. Blume-Kohout, S. Croke, and D. Gottesman, ``Streaming universal distortion-free entanglement concentration,'' 2009. \href{https://arxiv.org/abs/0910.5952}{arXiv:0910.5952}.
\bibitem{shevtsova} I. Shevtsova, ``On the absolute constants in the Berry--Esseen type inequalities for identically distributed summands,'' 2011. \href{https://arxiv.org/abs/1111.6554}{arXiv:1111.6554}.
\bibitem{barenco} A. Barenco \emph{et al.}, ``Elementary gates for quantum computation,'' \emph{Phys. Rev. A}, vol.~52, pp.~3457--3467, 1995. \href{https://arxiv.org/abs/quant-ph/9503016}{arXiv:quant-ph/9503016}.
\bibitem{shende} V.~V. Shende, S.~S. Bullock, and I.~L. Markov, ``Synthesis of quantum-logic circuits,'' \emph{IEEE Trans. Comput.-Aided Des. Integr. Circuits Syst.}, vol.~25, pp.~1000--1010, 2006. \href{https://arxiv.org/abs/quant-ph/0406176}{arXiv:quant-ph/0406176}.
\bibitem{dawson} C.~M. Dawson and M.~A. Nielsen, ``The Solovay--Kitaev algorithm,'' \emph{Quantum Inf. Comput.}, vol.~6, pp.~81--95, 2006. \href{https://arxiv.org/abs/quant-ph/0505030}{arXiv:quant-ph/0505030}.
\bibitem{bch} D. Bacon, I.~L. Chuang, and A.~W. Harrow, ``The quantum Schur transform: I. Efficient qudit circuits,'' 2006. \href{https://arxiv.org/abs/quant-ph/0601001}{arXiv:quant-ph/0601001}.
\bibitem{kirby} W.~M. Kirby and F.~W. Strauch, ``A practical quantum algorithm for the Schur transform,'' 2017. \href{https://arxiv.org/abs/1709.07119}{arXiv:1709.07119}.
\bibitem{krovi} H. Krovi, ``An efficient high dimensional quantum Schur transform,'' \emph{Quantum}, vol.~3, p.~122, 2019. \href{https://doi.org/10.22331/q-2019-02-14-122}{doi:10.22331/q-2019-02-14-122}.
\bibitem{kondra} T.~V. Kondra, P. Barrios Hita, J. Neumann, H. Kampermann, and D. Bru\ss, ``Fundamental limitations on entanglement extraction from purity,'' 2026. \href{https://arxiv.org/abs/2605.29197}{arXiv:2605.29197}.
\bibitem{zhang-extraction} Y. Zhang, ``Global-unitary entanglement extraction: capacity and exact strong-converse exponents,'' repository preprint, 2026. \href{https://github.com/mxym/math/tree/44d096fe63e4b1ef9462a338b941a11ed931b790/research/global-unitary-ppt-extraction}{Fixed repository source}.
\bibitem{zhang-universal} Y. Zhang, ``State-independent entanglement extraction: exact exponents, reliability and a Gaussian threshold,'' repository preprint, 2026. \href{https://github.com/mxym/math/tree/44d096fe63e4b1ef9462a338b941a11ed931b790/research/universal-schur-extraction}{Fixed repository source}.
\bibitem{zhang-appt-entropy} Y. Zhang, ``A relative APPT entropy equivalent uniform in both dimensions and every order,'' repository manuscript, 2026. \href{https://github.com/mxym/math/tree/44d096fe63e4b1ef9462a338b941a11ed931b790/research/appt-all-dimension-entropy-profile-20261015}{Fixed repository source}.
\bibitem{zhang-appt-purity} Y. Zhang, ``APPT purity: qutrit exact theorems and high-dimensional asymptotics,'' repository manuscript and verified sources, 2026. \href{https://github.com/mxym/math/tree/44d096fe63e4b1ef9462a338b941a11ed931b790/formalizations/appt-high-dimensional-review}{Fixed formalization and scope records}.
\end{thebibliography}

\end{document}
